% Shared preamble for the Computer Science course notes.
% Each course has a main.tex that does % Shared preamble for the Computer Science course notes.
% Each course has a main.tex that does % Shared preamble for the Computer Science course notes.
% Each course has a main.tex that does % Shared preamble for the Computer Science course notes.
% Each course has a main.tex that does \input{../common/preamble} and then
% \include's one file per chapter/lecture from its chapters/ directory.
%------------------------------------------------------------------------------
\usepackage[utf8]{inputenc}
\usepackage[T1]{fontenc}
\usepackage[english]{babel}
\usepackage{lmodern}
\usepackage[margin=1in]{geometry}
\usepackage{mathtools}
\usepackage{amsmath}
\usepackage{amssymb}
\usepackage{amsthm}
\usepackage{enumitem}
\usepackage{booktabs}
\usepackage{array}
\usepackage{graphicx}
\usepackage{float}
\usepackage{xcolor}
\usepackage{listings}
\usepackage{tikz}
\usepackage[most]{tcolorbox}
\usepackage[hidelinks]{hyperref}

\usetikzlibrary{arrows.meta,positioning,shapes,calc,trees,automata}

%------------------------------------------------------------------------------
% Math shortcuts (same names as the Math 239 notes)
\newcommand{\N}{\mathbb{N}}
\newcommand{\Z}{\mathbb{Z}}
\newcommand{\Q}{\mathbb{Q}}
\newcommand{\R}{\mathbb{R}}
\newcommand{\C}{\mathbb{C}}
\newcommand{\bigO}{\mathcal{O}}
\newcommand{\floor}[1]{\left\lfloor #1 \right\rfloor}
\newcommand{\ceil}[1]{\left\lceil #1 \right\rceil}
\newcommand{\abs}[1]{\left\lvert #1 \right\rvert}
\newcommand{\code}[1]{\texttt{#1}}

%------------------------------------------------------------------------------
% Theorem-like environments (numbered per chapter)
\theoremstyle{definition}
\newtheorem{definition}{Definition}[chapter]
\newtheorem{example}[definition]{Example}
\theoremstyle{plain}
\newtheorem{theorem}[definition]{Theorem}
\newtheorem{lemma}[definition]{Lemma}
\newtheorem{corollary}[definition]{Corollary}
\theoremstyle{remark}
\newtheorem*{remark}{Remark}

% Boxes: "important" for exam-relevant facts, "aside" for side comments.
\newtcolorbox{important}[1][]{colback=blue!4,colframe=blue!50!black,
  fonttitle=\bfseries,title={#1},boxrule=1pt}
\newtcolorbox{aside}{colback=gray!5,colframe=gray!60,boxrule=0.4pt}

%------------------------------------------------------------------------------
% Code listings
\definecolor{codebg}{gray}{0.96}
\definecolor{codecomment}{rgb}{0.25,0.5,0.35}
\definecolor{codekeyword}{rgb}{0.1,0.2,0.6}
\definecolor{codestring}{rgb}{0.6,0.1,0.1}
\lstset{
  basicstyle=\ttfamily\small,
  backgroundcolor=\color{codebg},
  commentstyle=\color{codecomment}\itshape,
  keywordstyle=\color{codekeyword}\bfseries,
  stringstyle=\color{codestring},
  columns=fullflexible,
  keepspaces=true,
  showstringspaces=false,
  breaklines=true,
  frame=single,
  rulecolor=\color{gray!40},
  tabsize=4,
  upquote=true,
  literate={~}{{\textasciitilde}}1
}

\lstdefinelanguage{Racket}{
  morekeywords={define,lambda,cond,else,if,let,let*,letrec,local,struct,
    define-struct,and,or,not,empty,cons,first,rest,list,check-expect,
    require,provide,begin,set!,when,unless,quote},
  sensitive=true,
  morecomment=[l]{;},
  morecomment=[s]{\#|}{|\#},
  morestring=[b]",
  alsoletter={-?!*>/+=<},
}

\lstdefinelanguage{MIPS}{
  morekeywords={add,addi,addu,sub,subu,mult,multu,div,divu,mfhi,mflo,lis,
    lw,sw,slt,sltu,beq,bne,jr,jalr,j,jal,and,or,xor,nor,sll,srl,lui,ori,word},
  sensitive=false,
  morecomment=[l]{;},
  morecomment=[l]{\#},
  alsoletter={.},
}

\lstdefinestyle{c}{language=C}
\lstdefinestyle{cpp}{language=C++}
\lstdefinestyle{racket}{language=Racket}
\lstdefinestyle{bash}{language=bash}
\lstdefinestyle{mips}{language=MIPS}
 and then
% \include's one file per chapter/lecture from its chapters/ directory.
%------------------------------------------------------------------------------
\usepackage[utf8]{inputenc}
\usepackage[T1]{fontenc}
\usepackage[english]{babel}
\usepackage{lmodern}
\usepackage[margin=1in]{geometry}
\usepackage{mathtools}
\usepackage{amsmath}
\usepackage{amssymb}
\usepackage{amsthm}
\usepackage{enumitem}
\usepackage{booktabs}
\usepackage{array}
\usepackage{graphicx}
\usepackage{float}
\usepackage{xcolor}
\usepackage{listings}
\usepackage{tikz}
\usepackage[most]{tcolorbox}
\usepackage[hidelinks]{hyperref}

\usetikzlibrary{arrows.meta,positioning,shapes,calc,trees,automata}

%------------------------------------------------------------------------------
% Math shortcuts (same names as the Math 239 notes)
\newcommand{\N}{\mathbb{N}}
\newcommand{\Z}{\mathbb{Z}}
\newcommand{\Q}{\mathbb{Q}}
\newcommand{\R}{\mathbb{R}}
\newcommand{\C}{\mathbb{C}}
\newcommand{\bigO}{\mathcal{O}}
\newcommand{\floor}[1]{\left\lfloor #1 \right\rfloor}
\newcommand{\ceil}[1]{\left\lceil #1 \right\rceil}
\newcommand{\abs}[1]{\left\lvert #1 \right\rvert}
\newcommand{\code}[1]{\texttt{#1}}

%------------------------------------------------------------------------------
% Theorem-like environments (numbered per chapter)
\theoremstyle{definition}
\newtheorem{definition}{Definition}[chapter]
\newtheorem{example}[definition]{Example}
\theoremstyle{plain}
\newtheorem{theorem}[definition]{Theorem}
\newtheorem{lemma}[definition]{Lemma}
\newtheorem{corollary}[definition]{Corollary}
\theoremstyle{remark}
\newtheorem*{remark}{Remark}

% Boxes: "important" for exam-relevant facts, "aside" for side comments.
\newtcolorbox{important}[1][]{colback=blue!4,colframe=blue!50!black,
  fonttitle=\bfseries,title={#1},boxrule=1pt}
\newtcolorbox{aside}{colback=gray!5,colframe=gray!60,boxrule=0.4pt}

%------------------------------------------------------------------------------
% Code listings
\definecolor{codebg}{gray}{0.96}
\definecolor{codecomment}{rgb}{0.25,0.5,0.35}
\definecolor{codekeyword}{rgb}{0.1,0.2,0.6}
\definecolor{codestring}{rgb}{0.6,0.1,0.1}
\lstset{
  basicstyle=\ttfamily\small,
  backgroundcolor=\color{codebg},
  commentstyle=\color{codecomment}\itshape,
  keywordstyle=\color{codekeyword}\bfseries,
  stringstyle=\color{codestring},
  columns=fullflexible,
  keepspaces=true,
  showstringspaces=false,
  breaklines=true,
  frame=single,
  rulecolor=\color{gray!40},
  tabsize=4,
  upquote=true,
  literate={~}{{\textasciitilde}}1
}

\lstdefinelanguage{Racket}{
  morekeywords={define,lambda,cond,else,if,let,let*,letrec,local,struct,
    define-struct,and,or,not,empty,cons,first,rest,list,check-expect,
    require,provide,begin,set!,when,unless,quote},
  sensitive=true,
  morecomment=[l]{;},
  morecomment=[s]{\#|}{|\#},
  morestring=[b]",
  alsoletter={-?!*>/+=<},
}

\lstdefinelanguage{MIPS}{
  morekeywords={add,addi,addu,sub,subu,mult,multu,div,divu,mfhi,mflo,lis,
    lw,sw,slt,sltu,beq,bne,jr,jalr,j,jal,and,or,xor,nor,sll,srl,lui,ori,word},
  sensitive=false,
  morecomment=[l]{;},
  morecomment=[l]{\#},
  alsoletter={.},
}

\lstdefinestyle{c}{language=C}
\lstdefinestyle{cpp}{language=C++}
\lstdefinestyle{racket}{language=Racket}
\lstdefinestyle{bash}{language=bash}
\lstdefinestyle{mips}{language=MIPS}
 and then
% \include's one file per chapter/lecture from its chapters/ directory.
%------------------------------------------------------------------------------
\usepackage[utf8]{inputenc}
\usepackage[T1]{fontenc}
\usepackage[english]{babel}
\usepackage{lmodern}
\usepackage[margin=1in]{geometry}
\usepackage{mathtools}
\usepackage{amsmath}
\usepackage{amssymb}
\usepackage{amsthm}
\usepackage{enumitem}
\usepackage{booktabs}
\usepackage{array}
\usepackage{graphicx}
\usepackage{float}
\usepackage{xcolor}
\usepackage{listings}
\usepackage{tikz}
\usepackage[most]{tcolorbox}
\usepackage[hidelinks]{hyperref}

\usetikzlibrary{arrows.meta,positioning,shapes,calc,trees,automata}

%------------------------------------------------------------------------------
% Math shortcuts (same names as the Math 239 notes)
\newcommand{\N}{\mathbb{N}}
\newcommand{\Z}{\mathbb{Z}}
\newcommand{\Q}{\mathbb{Q}}
\newcommand{\R}{\mathbb{R}}
\newcommand{\C}{\mathbb{C}}
\newcommand{\bigO}{\mathcal{O}}
\newcommand{\floor}[1]{\left\lfloor #1 \right\rfloor}
\newcommand{\ceil}[1]{\left\lceil #1 \right\rceil}
\newcommand{\abs}[1]{\left\lvert #1 \right\rvert}
\newcommand{\code}[1]{\texttt{#1}}

%------------------------------------------------------------------------------
% Theorem-like environments (numbered per chapter)
\theoremstyle{definition}
\newtheorem{definition}{Definition}[chapter]
\newtheorem{example}[definition]{Example}
\theoremstyle{plain}
\newtheorem{theorem}[definition]{Theorem}
\newtheorem{lemma}[definition]{Lemma}
\newtheorem{corollary}[definition]{Corollary}
\theoremstyle{remark}
\newtheorem*{remark}{Remark}

% Boxes: "important" for exam-relevant facts, "aside" for side comments.
\newtcolorbox{important}[1][]{colback=blue!4,colframe=blue!50!black,
  fonttitle=\bfseries,title={#1},boxrule=1pt}
\newtcolorbox{aside}{colback=gray!5,colframe=gray!60,boxrule=0.4pt}

%------------------------------------------------------------------------------
% Code listings
\definecolor{codebg}{gray}{0.96}
\definecolor{codecomment}{rgb}{0.25,0.5,0.35}
\definecolor{codekeyword}{rgb}{0.1,0.2,0.6}
\definecolor{codestring}{rgb}{0.6,0.1,0.1}
\lstset{
  basicstyle=\ttfamily\small,
  backgroundcolor=\color{codebg},
  commentstyle=\color{codecomment}\itshape,
  keywordstyle=\color{codekeyword}\bfseries,
  stringstyle=\color{codestring},
  columns=fullflexible,
  keepspaces=true,
  showstringspaces=false,
  breaklines=true,
  frame=single,
  rulecolor=\color{gray!40},
  tabsize=4,
  upquote=true,
  literate={~}{{\textasciitilde}}1
}

\lstdefinelanguage{Racket}{
  morekeywords={define,lambda,cond,else,if,let,let*,letrec,local,struct,
    define-struct,and,or,not,empty,cons,first,rest,list,check-expect,
    require,provide,begin,set!,when,unless,quote},
  sensitive=true,
  morecomment=[l]{;},
  morecomment=[s]{\#|}{|\#},
  morestring=[b]",
  alsoletter={-?!*>/+=<},
}

\lstdefinelanguage{MIPS}{
  morekeywords={add,addi,addu,sub,subu,mult,multu,div,divu,mfhi,mflo,lis,
    lw,sw,slt,sltu,beq,bne,jr,jalr,j,jal,and,or,xor,nor,sll,srl,lui,ori,word},
  sensitive=false,
  morecomment=[l]{;},
  morecomment=[l]{\#},
  alsoletter={.},
}

\lstdefinestyle{c}{language=C}
\lstdefinestyle{cpp}{language=C++}
\lstdefinestyle{racket}{language=Racket}
\lstdefinestyle{bash}{language=bash}
\lstdefinestyle{mips}{language=MIPS}
 and then
% \include's one file per chapter/lecture from its chapters/ directory.
%------------------------------------------------------------------------------
\usepackage[utf8]{inputenc}
\usepackage[T1]{fontenc}
\usepackage[english]{babel}
\usepackage{lmodern}
\usepackage[margin=1in]{geometry}
\usepackage{mathtools}
\usepackage{amsmath}
\usepackage{amssymb}
\usepackage{amsthm}
\usepackage{enumitem}
\usepackage{booktabs}
\usepackage{array}
\usepackage{graphicx}
\usepackage{float}
\usepackage{xcolor}
\usepackage{listings}
\usepackage{tikz}
\usepackage[most]{tcolorbox}
\usepackage[hidelinks]{hyperref}

\usetikzlibrary{arrows.meta,positioning,shapes,calc,trees,automata}

%------------------------------------------------------------------------------
% Math shortcuts (same names as the Math 239 notes)
\newcommand{\N}{\mathbb{N}}
\newcommand{\Z}{\mathbb{Z}}
\newcommand{\Q}{\mathbb{Q}}
\newcommand{\R}{\mathbb{R}}
\newcommand{\C}{\mathbb{C}}
\newcommand{\bigO}{\mathcal{O}}
\newcommand{\floor}[1]{\left\lfloor #1 \right\rfloor}
\newcommand{\ceil}[1]{\left\lceil #1 \right\rceil}
\newcommand{\abs}[1]{\left\lvert #1 \right\rvert}
\newcommand{\code}[1]{\texttt{#1}}

%------------------------------------------------------------------------------
% Theorem-like environments (numbered per chapter)
\theoremstyle{definition}
\newtheorem{definition}{Definition}[chapter]
\newtheorem{example}[definition]{Example}
\theoremstyle{plain}
\newtheorem{theorem}[definition]{Theorem}
\newtheorem{lemma}[definition]{Lemma}
\newtheorem{corollary}[definition]{Corollary}
\theoremstyle{remark}
\newtheorem*{remark}{Remark}

% Boxes: "important" for exam-relevant facts, "aside" for side comments.
\newtcolorbox{important}[1][]{colback=blue!4,colframe=blue!50!black,
  fonttitle=\bfseries,title={#1},boxrule=1pt}
\newtcolorbox{aside}{colback=gray!5,colframe=gray!60,boxrule=0.4pt}

%------------------------------------------------------------------------------
% Code listings
\definecolor{codebg}{gray}{0.96}
\definecolor{codecomment}{rgb}{0.25,0.5,0.35}
\definecolor{codekeyword}{rgb}{0.1,0.2,0.6}
\definecolor{codestring}{rgb}{0.6,0.1,0.1}
\lstset{
  basicstyle=\ttfamily\small,
  backgroundcolor=\color{codebg},
  commentstyle=\color{codecomment}\itshape,
  keywordstyle=\color{codekeyword}\bfseries,
  stringstyle=\color{codestring},
  columns=fullflexible,
  keepspaces=true,
  showstringspaces=false,
  breaklines=true,
  frame=single,
  rulecolor=\color{gray!40},
  tabsize=4,
  upquote=true,
  literate={~}{{\textasciitilde}}1
}

\lstdefinelanguage{Racket}{
  morekeywords={define,lambda,cond,else,if,let,let*,letrec,local,struct,
    define-struct,and,or,not,empty,cons,first,rest,list,check-expect,
    require,provide,begin,set!,when,unless,quote},
  sensitive=true,
  morecomment=[l]{;},
  morecomment=[s]{\#|}{|\#},
  morestring=[b]",
  alsoletter={-?!*>/+=<},
}

\lstdefinelanguage{MIPS}{
  morekeywords={add,addi,addu,sub,subu,mult,multu,div,divu,mfhi,mflo,lis,
    lw,sw,slt,sltu,beq,bne,jr,jalr,j,jal,and,or,xor,nor,sll,srl,lui,ori,word},
  sensitive=false,
  morecomment=[l]{;},
  morecomment=[l]{\#},
  alsoletter={.},
}

\lstdefinestyle{c}{language=C}
\lstdefinestyle{cpp}{language=C++}
\lstdefinestyle{racket}{language=Racket}
\lstdefinestyle{bash}{language=bash}
\lstdefinestyle{mips}{language=MIPS}
